\documentclass{beamer}

% TODO (GM): Styling!
\usetheme{Copenhagen}
\setbeamertemplate{headline}{}

\usepackage{tikz}
\usepackage{pgfplots}
\usepackage{amsmath}
\usepackage{mathtools}
\usepackage[ngerman]{babel}

\usepackage{relsize}
\usepackage{algorithm}
\usepackage{algorithmic}
\usepackage{xcolor}
\usepackage{todonotes}

\presetkeys{todonotes}{inline}{}

\usetikzlibrary{overlay-beamer-styles}
\usetikzlibrary{arrows.meta,calc}

% -- BEGIN SECTION TOKENS AND CUSTOM COMMANDS --

\newcommand{\sectionstart}[1]{
    \section{#1}
    \begin{frame}
        \centering
        \huge{#1}
    \end{frame}
}

\newcommand{\R}{\mathbb{R}}
\newcommand{\E}{\mathbb{E}}
% -- END SECTION TOKENS AND CUSTOM COMMANDS --

% -- BEGIN SECTION SLIDES --

\title[]{Brownian Particle Systems and its\\Mean-Field Limit}
\author{Max Grübmeyer}
\date{17th of July, 2026}

\begin{document}
    \begin{frame}
        \titlepage
    \end{frame}

    \begin{frame}{Brownian Particle Systems}
        Consider a particle system in $\R^d$ given by the SDE
        \begin{equation*}
            X_i(t) = X_i(0) + \sqrt2B_i(t)
        \end{equation*}

        \pause
        
        and the particle system given by
        \begin{equation*}
            dX_i(t) = \sqrt2dB_i(t) + b(X_i(t),\mu_t^N)dt
        \end{equation*}
    \end{frame}

    \sectionstart{A Simple Start}

    \begin{frame}{Problem formulation}
        Consider a particle system in $\R^d$ given by the SDE
        \begin{equation*}
            X_i(t) = X_i(0) + \sqrt2B_i(t)
        \end{equation*}
        for $1\leq i\leq N$ and $B_i$ being independent standard Brownian motions
    \end{frame}

    \begin{frame}{Empirical Measures}
        Consider the empirical measure $\mu_t^N$ given by
        \begin{equation*}
            \mu_t^N = \frac{1}{N}\sum_{i=1}^N\delta_{X_i(t)}
        \end{equation*}

        \pause

        The canonical duality pairing is therefore given by
        \begin{equation*}
            \langle\mu_t^N,\varphi\rangle = \frac{1}{N}\sum_{i=1}^N\varphi(X_i(t))
        \end{equation*}
        for a test function $\varphi\in\mathcal{D}(\R^d)$
    \end{frame}

    \sectionstart{Formal evolution and law of large numbers}

    \begin{frame}{Formal Evolution}
        Take $\varphi(X_i(t))$ and apply Itô's formula to obtain
        \begin{align*}
            d\varphi(X(t))
            %&= \sqrt2\sum_{j=1}^d\partial_{x_j}\varphi(X(t))dB^j(s)
            %+ \sum_{j,l=1}^d\partial_{x_i,x_j}^2\varphi(X(t))dB^j(s)dB^l(s)\\
            &= \sqrt2\nabla\varphi(X(s))\cdot dB(s) + \Delta\varphi(X(s))ds
        \end{align*}

        \pause
        Integrating w.r.t. time and averaging over all particles yields
        \begin{align*}
            \frac{1}{N}\sum_{i=1}^N\varphi(X_i(t))
            &= \frac{1}{N}\sum_{i=1}^N\varphi(X_i(0))
            + \frac{\sqrt2}{N}\sum_{i=1}^N\int_0^t\nabla\varphi(X_i(s))\cdot dB_i(s)\\
            &+ \frac{1}{N}\sum_{i=1}^N\int_0^t\Delta\varphi(X_i(s))ds
        \end{align*}
    \end{frame}

    \begin{frame}{Formal Evolution}
        Define
        \begin{equation*}
            M_t^N(\varphi)
            \coloneq
            \frac{\sqrt 2}{N}\sum_{i=1}^N\int_0^t\nabla\varphi\cdot dB_i(s)
        \end{equation*}
    
        \pause
        Then the previous equality can be rewritten as
        \begin{equation*}
            \langle\mu_t^N, \varphi\rangle = \langle\mu_0^N, \varphi\rangle
            +\int_0^t\langle\mu_s^N,\Delta\varphi\rangle ds + M_t^N
        \end{equation*}
    \end{frame}

    \begin{frame}{Second moment of $M_t^N$}
        When looking at the second moment we obtain
        \begin{align*}
            \E[M_t^N(\varphi)^2]
            &= \E\left[\left(
            \frac{\sqrt 2}{N}\sum_{i=1}^N\int_0^t\nabla\varphi\cdot dB_i(s)
            \right)^2\right]\\
            &= \frac{2}{N^2}\E\left[
            \sum_{i=1}^N\int_0^t(\nabla\varphi(X_i(s)\cdot dB_i(s))^2\right]\\
            &= \frac{2}{N}\E\left[
            \int_0^t(\nabla\varphi(X_1(s)\cdot dB_1(s))^2\right]\\
            &= \frac{2}{N}\E\left[\int_0^t\Vert\nabla\varphi(X_i(s))\Vert_2^2ds\right]\\
            &\leq\frac{2T}{N}\Vert\nabla\varphi\Vert_\infty^2
        \end{align*}
    \end{frame}

    \begin{frame}{Weak Convergence}
        Take a look at the decomposition again
        \begin{equation*}
            \langle\mu_t^N, \varphi\rangle = \langle\mu_0^N, \varphi\rangle
            +\int_0^t\langle\mu_s^N,\Delta\varphi\rangle ds + M_t^N(\varphi)
        \end{equation*}

        \pause
        \begin{itemize}
            \item $M_t^N(\varphi)$: vanishes in $L^2$
            \pause
            \item $\langle\mu_0^N,\varphi\rangle$ converges by SLLN
            \pause
            \item $\int_0^t\langle\mu_s^N,\varphi\rangle ds$ converges by SLLN + DCT
        \end{itemize}
    \end{frame}

    \sectionstart{Limit equation}

    \begin{frame}{Weak formulation}
        Thus the whole measure (for a fixed time $t$) converges to
        \begin{equation*}
            \langle\rho(t,x)dx,\varphi\rangle = \langle\rho(0,x)dx,\varphi\rangle
            + \int_0^t\langle\rho(s,x)dx,\Delta\varphi\rangle ds
        \end{equation*}

        \pause
        
        and we obtain
        \begin{equation*}
            \frac{d}{dt}\langle\rho(t,x)dx,\varphi\rangle
            = \langle\partial_t\rho(t,x)dx,\varphi\rangle
            = \langle\rho(t,x)dx,\Delta\varphi\rangle
        \end{equation*}
    \end{frame}

    \begin{frame}{A more explicit form}
        Since $\text{Var}(X_i(t)) = \text{Var}(\sqrt 2B_i(t)) = 2\text{Var}(B_i(t))$
        the transition density of our particles is given by
        \begin{equation*}
            p_t(x,y) = \frac{1}{\sqrt{4\pi t}^d}\exp\left(-\frac{\vert x-y\vert^2}{4t}\right)
        \end{equation*}

        \pause

        Recall the heat kernel for $t > 0$ is given by
        \begin{equation*}
            \Phi(x,t) =\frac{1}{\sqrt{4\pi t}^d}\exp\left(-\frac{\vert x\vert^2}{4t}\right)
        \end{equation*}

        \pause

        Thus by the law of total probability we have
        \begin{equation*}
            \rho(t,x)dx = \int_{\R^d}p_t(x,y)\rho_0(y)dy = (\Phi(\cdot,t)\ast\rho_0)(x)
        \end{equation*}
    \end{frame}

    \sectionstart{A more general problem}

    \begin{frame}{Problem Introduction}
        Consider the particle system given by the equations
        \begin{equation*}
            dX_i^N(t) =  \sqrt2B_i(t) + b(X_i^N(t),\mu_t^N)dt
        \end{equation*}
        for $1\leq i\leq N$ where $B_i$ are independent standard Brownian 
        motions.

        \pause
        \medskip

        Where $b$ has the form
        \begin{equation*}
            b(X_i^N(t),\mu_t^N) = \int_{\R^d}\beta(X_i^N(t),y)\,d\mu_t^N(y)
        \end{equation*}
        for $\beta$ Lipschitz and bounded.
    \end{frame}

    \begin{frame}{Limit Candidate}
        Consider the nonlinear equation
        \begin{equation*}
            d\bar X_i(t) = \sqrt 2 dB_i(t) + b(\bar X_i(t),\rho_t)dt
        \end{equation*}
        where $\rho_t$ is the law of $\bar X(t)$
    \end{frame}

    \sectionstart{Existence and uniqueness of the nonlinear process}

    \begin{frame}{Wasserstein Metric}
        Let $C = C([0,T],\R^d)$ and let $\mathcal{M} = \mathcal{M}(C)$
        be the space of all probability measures on $C$.

        \pause
        \medskip
    
        \begin{block}{The Wasserstein metric}
            Given two probability measures $\mu,\nu\in \mathcal{M}$
            the (capped) Wasserstein metric
            is defined as
            \begin{equation*}
                D_T(\mu,\nu)
                = \inf_{\gamma\in\Gamma(\mu,\nu)}\int_{\mathcal{M}\times\mathcal{M}}
                \left(\sup_{s\leq T}\vert\omega_1(s)-\omega_2(s)\vert \wedge 1\right)
                d\gamma(\omega_1,\omega_2)
            \end{equation*}
            defines a complete metric on $\mathcal{M}$.
        \end{block}
    \end{frame}

    \begin{frame}{Proof Sketch: Unique Existence via Contraction}
        Define a map $\Phi: \mathcal{M}(C) \to \mathcal{M}(C)$ that assigns to a measure $m$ the law of the solution to the SDE:
        \begin{equation*}
            dX(t) = \sqrt 2dB(t) + b(X(t), m_t)dt
        \end{equation*}

        \pause

        Notice that $m$ is a FP of $\Phi$ if and only if the associated process solves
        the nonlinear equation.
    \end{frame}
    \begin{frame}{Prook Sketch: Unique Existence via Contraction}
        By Gronwall, we have the following estimate
        \begin{equation*}
            D_T(\Phi(m_1), \Phi(m_2)) \leq C_T \int_0^T D_u(m_1, m_2) du
        \end{equation*}
        \pause

        Iterate this estimate to obtain
        \begin{equation*}
            D_T(\Phi^{k+1}(m),\Phi^k(m)) \leq C_T \frac{T^k}{k!}D_T(\Phi(m),m)
        \end{equation*}
    \end{frame}

    \sectionstart{The central convergence result}

    \begin{frame}{The central result}
        \begin{block}{The central convergence result}
            For each $T>0$ we have the following estimate
            \begin{equation*}
                \sup_N\sqrt N\E\left[\sup_{t\leq T}\vert X_i^N(t)-\bar X_i(t)\vert\right] < 
                \infty
            \end{equation*}
        \end{block}
    \end{frame}
    
    \begin{frame}{Proof Sketch: Setup}
        Bound the deviation between the particle system $X_i^N(t)$
        and the mean-field limit $\bar{X}_i(t)$.

        \pause
        \medskip
        
        Define the centered interaction term to isolate fluctuations:
        \begin{equation*}
            b_s(x, x') := b(x, x') - \int b(x, y) \mu_s(dy)
        \end{equation*}

        \pause

        In particular note that for a known $x$ we have
        \begin{equation*}
            \E[b_s(x,\bar X_j(s)] = 0
        \end{equation*}
    \end{frame}
    
    \begin{frame}{Proof Sketch: Symmetry Gronwall Inequality}
        We get the following bound by Lipschitz property of the drift, symmetry,
        and Gronwall
        \begin{equation*}
            N\mathbb{E}[|X_i^N(t) - \bar{X}_i(t)|] \leq K_T \int_0^t \sum_{i=1}^N \mathbb{E} \left[ \left| \frac{1}{N} \sum_{j=1}^N b_s(\bar X_i(s), \bar X_j(s)) \right| \right] ds
        \end{equation*}
    \end{frame}
    
    \begin{frame}{Proof Sketch: Variance Analysis}
        Now the second moment of the interaction term:
        \begin{align*}
            \mathbb{E} \left[ \left| \frac{1}{N} \sum_{j=1}^N b_s(\bar X_i, \bar X_j)
            \right|^2 \right]
            &= \frac{1}{N^2} \sum_{j,k=1}^N
            \mathbb{E}[b_s(\bar X_i, \bar X_j) b_s(\bar X_i, \bar X_k)]\\
            &=\frac{1}{N^2}\sum_{j=1}^N\E[b_s(\bar X_i, \bar X_j)b_s(\bar X_i, \bar X_j)]
        \end{align*}

        since $\mathbb{E}[b_s | \bar X_i] = 0$.

        \pause
        \bigskip

        Thus the sum yields an $O(1/N)$ bound.
        Since $\E\vert Z\vert \leq \sqrt{\E[Z^2]}$ by Jensen this completes the proof.
    \end{frame}

    \section{Weak Convergence of Measures}

    \begin{frame}{Weak Convergence of Measures}
        Let
        \begin{align*}
            A_N&\coloneq\frac{1}{N}\sum_{i=1}^N\varphi(X_i^N(t))-\varphi(\bar X_i(t))\\
            B_N&\coloneq\frac{1}{N}\sum_{i=1}^N\varphi(\bar X_i(t))-\langle\rho_t,\varphi\rangle
        \end{align*}

        \pause
    
        Take a look at the following estimate
        \begin{align*}
            &\vert\langle\mu_t^N-\rho_t,\varphi\rangle\vert
            \leq \vert A_N\vert + \vert B_N\vert            
        \end{align*}
        to achieve convergence.
    \end{frame}

    \sectionstart{Limit equation}

    \begin{frame}{Limit equation}
    Let $X(t)$ solve the nonlinear limiting equation
    \begin{equation*}
        dX(t)
        =
        \sqrt{2}\,dB(t)
        +
        b(X(t),\rho_t)\,dt ,
    \end{equation*}
    where
    \begin{equation*}
        \rho_t = \text{Law}(X(t)).
    \end{equation*}
\end{frame}


\begin{frame}{Evolution}
    The measure $\rho_t$ is characterized by its action on test functions:
    \begin{equation*}
        \langle \rho_t,\varphi\rangle
        =
        \int_{\mathbb{R}^d}\varphi(x)\rho_t(dx)
        =
        \mathbb{E}[\varphi(X(t))].
    \end{equation*}

    \pause
    \medskip
    
    For $\varphi\in \mathcal{D}(\mathbb{R}^d)$, Itô's formula gives
    \begin{equation*}
        d\varphi(X(t))
        =
        M_t(\varphi)+
        \Delta\varphi(X(t))dt
        +
        \nabla\varphi(X(t))
        \cdot
        b(X(t),\rho_t)\,dt .
    \end{equation*}
\end{frame}

\begin{frame}{Weak formulation}
    Plugging this in we obtain
    \begin{align*}
        \langle\rho_t,\varphi\rangle
        &=
        \langle\rho_0,\varphi\rangle
        +
        \int_0^t
        \left\langle
        \rho_s,
        \Delta\varphi
        \right\rangle ds\\
        &+
        \int_0^t
        \int_{\mathbb{R}^d}
        \nabla\varphi(x)
        \cdot
        \beta(x,\rho_s)
        \rho_s(dx)\,ds .
    \end{align*}

    \pause
    This is the weak form of the nonlinear Fokker--Planck equation
    \begin{equation*}
        \partial_t\rho = \Delta \rho - \nabla\cdot(b(\cdot,\rho)\rho)
    \end{equation*}
\end{frame}

\begin{frame}
    \centering\huge{Thank you for your attention!}
\end{frame}
    
\end{document}

% -- END SECTION SLIDES --